\chapter{Lessons from the Prosperity Top Teams}
\label{ch:top_team_lessons}
\chaptermark{Lessons from the top teams}

This chapter distils patterns from 10+ public writeups across
Prosperity~1--3.  Rather than re-deriving strategies
(\cref{ch:algo_rounds}), we focus on the \emph{meta-game}: the
operational, analytical, and psychological patterns that separate
top-10 finishers from teams with similar ideas who finish hundreds
of places lower.

\begin{backgroundbox}[Prerequisites]
\begin{itemize}
  \item \Cref{ch:prosperity_setup}: competition mechanics, the
        Python trading API, product archetypes.
  \item \Cref{ch:algo_rounds}: strategy--theory mapping for each
        algorithmic round (market making, stat arb, options, location
        arbitrage, informed-trader exploitation).
  \item \Cref{ch:backtesting}: overfitting traps, walk-forward
        validation, and parameter-plateau selection.
\end{itemize}
\end{backgroundbox}

%======================================================================
\section{What every top team did the same}
\label{sec:top_same}
%======================================================================

Across every top-10 writeup we examined, five patterns recur
reliably as necessary (though not sufficient) conditions for a top
finish.  The set spans Frankfurt Hedgehogs (2nd, P3), CMU Physics
(7th, P3), Linear Utility (2nd, P2), Stanford Cardinal (P1), Alpha
Animals (9th, P3), jmerle (9th P2, 25th P3), and others.

Prosperity~3 had over 12,000 teams and a heavy-tailed PnL
distribution: the top~10 earn 5--10$\times$ the median, and 1st
to 100th can span 1M~SeaShells.  Small edges compound over five
rounds and 50{,}000 timesteps.

\subsection{Tooling from day one}
\label{sec:tooling_day1}

Every top-10 team built or installed core tooling \emph{before
the competition started}.  Three mandatory components:

\begin{enumerate}
  \item \textbf{Backtester.}  A forked jmerle or custom build.
        Linear Utility's P2 backtester replicated the matching
        engine and output logs in the competition website's exact
        format, enabling direct backtest-vs-live comparison.
        Frankfurt Hedgehogs used forked jmerle for systematic
        parameter sweeps and the official website for validating
        bot-interaction behaviour the local tool could not
        replicate.

  \item \textbf{Visualiser / dashboard.}  Frankfurt Hedgehogs'
        custom dashboard showed order-book levels, trade markers
        (maker/taker/own fills), PnL, position, and a normalisation
        feature displaying prices relative to a chosen indicator
        (e.g.\ the ``Wall Mid'').  CMU Physics used Jasper's
        community visualiser, but in Round~3 it pushed their AWS
        Lambda over the 100\,MB limit, crashing the algorithm and
        wiping rolling-window state.  They removed it from
        subsequent submissions.

  \item \textbf{Analysis notebooks.}  Every top team ran EDA in
        Jupyter before writing strategy code.  Frankfurt Hedgehogs
        prototyped via quick vectorised backtests in notebooks,
        graduating promising ideas to the full simulation.  Linear
        Utility ran parameter grid searches, optimising
        systematically rather than by hand.
\end{enumerate}

\begin{keyfactbox}[Tooling priority]
Build backtester + visualiser + analysis notebooks \emph{before}
the competition, using prior-year data.  Frankfurt Hedgehogs
credit weeks of pre-competition prep as a primary driver of their
top finishes across P2 and P3.
\end{keyfactbox}

\subsection{Study prior years exhaustively}
\label{sec:study_prior}

Top teams treat prior-year writeups as primary source material.
Frankfurt Hedgehogs discovered Olivia's daily-extrema pattern on
Squid Ink, one of P3's highest-value signals, because they had
studied Stanford Cardinal (P1) and Jasper (P2), both documenting
similar informed-bot behaviour.  They state: ``Anyone who
carefully analyzed historical Prosperity~2 data or public
write-ups \ldots\ could have anticipated similar behaviors and
prepared detection logic in advance.''

CMU Physics did the same for Round~4 (locational arbitrage): their
Macaron strategy resembled P2 Round~2, and Chris Roberts had
placed 3rd on that round, so they re-implemented with high
confidence.

Linear Utility's P2 writeup itself became a key source for the
Macaron hidden-taker exploit, which Frankfurt Hedgehogs note had
been ``discussed already'' by Jasper and Linear Utility.

Reading writeups is only the first step.  The full lesson is to
\emph{identify recurring product archetypes}, build skeletons for
each before the competition, and slot new products into the
matching skeleton as sample data arrives.

\begin{keyfactbox}[Archetype recycling across editions]
Archetypes appearing in every Prosperity edition (P1--P3):
\begin{itemize}
  \item A \textbf{fixed-value} product for tutorial market making
        (Amethysts $\to$ Rainforest Resin).
  \item A \textbf{slow random walk} product for drifting-value
        market making (Starfruit $\to$ Kelp).
  \item A \textbf{basket / ETF} product for statistical arbitrage
        (Gift Basket $\to$ Picnic Baskets).
  \item A \textbf{derivative / option} product
        (Coconut options $\to$ Volcanic Rock Vouchers).
  \item A \textbf{cross-exchange conversion} product
        (Orchids $\to$ Macarons).
  \item An \textbf{informed bot} trading at daily extrema
        (unnamed P1 bot $\to$ Olivia).
\end{itemize}
Teams with skeleton code per archetype wrote P3 strategy logic
within hours of data release; teams discovering patterns from
scratch spent days on analysis that could have been done in
advance.
\end{keyfactbox}

\subsection{Simplicity over sophistication}
\label{sec:simplicity}

Winning strategies are almost never complex ML pipelines.  They
are overwhelmingly \textbf{fair-value estimate $+$ threshold $+$
inventory management}.

Frankfurt Hedgehogs' Rainforest Resin strategy (their most
consistent PnL contributor, $\sim$39{,}000 SeaShells/round) had
three rules: (1)~take any order crossing fair value~10{,}000;
(2)~place passive quotes inside the best bid/ask;
(3)~flatten inventory at fair value when skewed.  ``No
sophisticated logic or aggressiveness was needed.''

Their basket stat-arb used a \emph{fixed threshold}: long below
$-T$, short above $+T$.  No MA crossovers, z-scores, or vol
scaling.  They warn: ``if you can't explain why a strategy should
work from first principles, then any `outperformance' in
historical data is probably noise.''

CMU Physics' basket strategy used z-scores, but with a hard-coded
historical mean and short rolling-window standard deviation, still
simpler than textbook implementations.  Their Round~5 ``YOLO
Croissants'' was simpler still: detect Olivia's signal, go max
long Croissants via baskets as leverage.

\begin{pitfallbox}[Complexity as a trap]
Frankfurt Hedgehogs tried a logistic regression on Macarons using
sunlight and sugar-price features.  Despite significant p-values
and $\sim$25k/day in backtests, they abandoned it: low confidence
out-of-sample, serialisation overhead slowed the trader, and it
clashed with their conversion-arbitrage logic.  The simple
arbitrage they deployed instead made 80{,}000--100{,}000 SeaShells
with minimal risk.  Takeaway: a strategy you can implement
correctly under time pressure beats a theoretically superior one
you cannot trust.
\end{pitfallbox}

\subsection{Iterate, do not over-engineer}
\label{sec:iterate}

Top teams submit early and learn from live results rather than
perfecting strategies in simulation.

CMU Physics lived this.  After Round~3 they dropped from 7th to
241st because of two bugs: Jasper's visualiser exceeding memory (crashing
Lambda, wiping rolling windows) and a quadratic IV fit that
stopped tracking true IV on submission day.  They iterated:
removed the visualiser, switched to a short rolling window of
mid IV (boosting backtested PnL from 80k to 200k/day), and in
Round~4 posted the highest single-round algo PnL of \emph{all}
12{,}000+ teams (447{,}251 SeaShells), climbing back to 8th.

Frankfurt Hedgehogs re-optimised all parameters after every round,
treating each as a new calibration opportunity.

\begin{keyfactbox}[The iteration principle]
Submit a working strategy on day~1 of each round.  Use days~2--3
for live-data analysis, re-calibration, and bug fixes.  A live
submission's information is worth more than another day of
backtesting.
\end{keyfactbox}

\subsection{Risk-aware position management}
\label{sec:risk_aware}

Top teams compute worst-case losses before deploying.  CMU Physics
estimated unhedged Volcanic Rock Vouchers could cost 40{,}000
SeaShells/step worst case, but average delta was closer to 160
units (not 400), giving a realistic worst case $\sim$16{,}000, less
than the 40{,}000 saved by avoiding hedging spread costs.  They
took the risk.  Frankfurt Hedgehogs computed a 95\% VaR of
$\sim$50{,}000 for their mean-reversion leg (25\% of their
200{,}000 lead) and kept the position as a relative hedge.

The pattern: quantify downside before accepting it, and size
positions so the worst plausible outcome does not eliminate you.

Frankfurt Hedgehogs pushed this into \emph{relative} risk
management in the final round.  With a 200{,}000 lead, they could
have dropped mean-reversion entirely.  Instead, reasoning that
close competitors were trading mean reversion aggressively (if
the market mean-reverted, those teams would gain while Frankfurt
stood still), they kept a moderate position, hedging
\emph{ranking} risk rather than chasing EV.  VaR of 50{,}000 was
25\% of their lead: tournament-aware risk management.

\begin{intuitionbox}
Risk management in a Prosperity round is a two-player game: you
versus the leaderboard.  Behind, you need variance to catch up;
ahead, match competitors' strategies so bad luck hurts them
equally.  Same logic as index-fund relative-performance mandates:
underperforming the benchmark is worse than absolute losses.
\end{intuitionbox}


%======================================================================
\section{What differentiated the top 10}
\label{sec:top_different}
%======================================================================

Among teams with solid tooling, prior-year study, and simple
strategies, the following differentiators separated top~10 from
top~100.

\subsection{Bot-pattern exploitation}
\label{sec:bot_patterns}

The highest-value P3 signal was detecting the bot Olivia, who
bought at daily lows and sold at daily highs on multiple products
(Squid Ink, Croissants, others).  Frankfurt Hedgehogs identified
her \emph{before} trader IDs were revealed in Round~5 by filtering
for quantity-15 trades and observing their placement at daily
extrema.  They then used her inferred position to bias basket
spread thresholds, a cross-product integration that compounded the
edge.

CMU Physics independently spotted ``suspicious trade quantity~15
patterns at highs/lows for Squid Ink and Croissants'' in Round~2
but judged it too late to productise.  In Round~5 with trader IDs
visible, they went all-in on Croissants in Olivia's first-trade
direction.

Detection methodology: track daily running min/max price.  A
trade at a daily extreme in the expected direction (buy at min,
sell at max) is a potential Olivia signal.  New extrema
invalidate earlier signals: a new low after a supposed buy kills
the prior flag.

Pre-Round~5, the only tell was quantity: Olivia traded 15 lots.
Frankfurt Hedgehogs' dashboard had a quantity filter that made
quantity-15 clustering at extrema visually obvious: purpose-built
tooling exposed a pattern the raw trade log hid.

\begin{keyfactbox}[Olivia detection: the P3 meta-signal]
In P3, Olivia bought 15 lots at daily lows and sold 15 at daily
highs on at least three products.  Detection was worth 8{,}000+
SeaShells/round on Squid Ink (Frankfurt Hedgehogs' estimate),
20{,}000+/round on Croissants, plus indirect basket-threshold
bias.  Missing Olivia cost $\sim$30{,}000--50{,}000 per round.
\end{keyfactbox}

\subsection{Cross-product data reuse}
\label{sec:cross_product}

Frankfurt Hedgehogs used Olivia's inferred Croissants position
to bias basket stat-arb thresholds.  Croissants were $\sim$50\%
of both baskets' value, so Olivia's Croissants direction predicted
the basket spread's drift.  They shifted a $\pm 50$ threshold
asymmetrically: if Olivia was short, long-entry moved to $-80$
and short-entry to $+20$.

CMU Physics went further in Round~5.  Reasoning that trading
baskets in Olivia's Croissants direction strictly dominated
stat-arb, they went long both baskets to hold 1{,}050 Croissants
(vs.\ the 250 direct limit), partially hedging Jams/Djembes
exposure.  This ``YOLO'' made up to 120{,}000 on its best day.

The risk calculus was principled, not reckless.  A bad day's
Croissants high--low range was $\sim$40 SeaShells, giving a PnL
lower bound of $40 \times 1{,}050 = 42{,}000$.  Basket-premium
risk had a 300-SeaShell/basket worst case, 45{,}000 total, but
premiums were stationary (95\%/90\% confidence for Baskets~1/2),
making the extreme move near-impossible.  Expected premium loss
$\approx 0$; expected Olivia gain large and positive.
``First-principles justification'' applied to position sizing.

\subsection{Parameter plateau selection}
\label{sec:param_plateau}

Frankfurt Hedgehogs visualised backtested PnL across a 2-D
parameter grid for their basket strategy.  Rather than the
maximiser, they picked parameters from ``consistent, flat regions
of good performance, reducing sensitivity to slight shifts and
avoiding overfit disasters.''  This is the plateau approach of
\cref{ch:backtesting}: centre of a plateau, not peak of a spike.

CMU Physics saw the flip side on Volcanic Rock z-score trading.
A colleague's strategy produced 150k/day in backtests, but
``small tweaks \ldots\ would lead to wildly different backtesting
results, some even heavily negative.''  They declined to deploy,
correctly applying the plateau principle at the cost of potential
PnL.

\subsection{Options: the IV scalping vs.\ mean-reversion debate}
\label{sec:options_debate}

Volcanic Rock Vouchers (Round~3) produced the largest PnL
dispersion among top teams.  The choices here show theory plus
empirical discipline.

Frankfurt Hedgehogs ran two alpha sources.  The primary was \emph{IV
scalping}: fit a parabolic smile to observed IVs across strikes,
price via Black--Scholes, trade deviations.  (In plain English:
implied volatility is what option prices \emph{imply} the market
thinks future volatility will be; scalping means trading short,
mean-reverting deviations from a smooth fitted curve.)  They validated by
testing for 1-lag negative autocorrelation in option returns
(short-term IV mean-reversion, necessary for scalping).
Secondary: lightweight EMA mean-reversion on underlying Volcanic
Rock with fixed thresholds.

Results: IV scalping delivered a stable 100{,}000--150{,}000/round;
mean reversion gave $+100{,}000$, $-50{,}000$, $-10{,}000$ across
three rounds.  The hybrid protected against missing a
mean-reversion regime while anchoring PnL to the reliable
scalping edge.

CMU Physics used the same quadratic smile but as \emph{sole} fair
value.  When the fit diverged from market IV on submission day,
they entered large directionals and held, finishing with only
75{,}755 SeaShells versus 200{,}000+ for the top.  Post-round: a short
rolling window of mid IV would have tracked the market far better,
boosting backtested PnL from 80k to 200k/day.

General principle: \emph{adaptive} models (rolling windows, EMAs)
beat \emph{static} fitted curves when the DGP (data-generating
process, the underlying statistical law producing the observed
prices) can shift between sample and evaluation.

\subsection{The ``Wall Mid'' and fair-value estimation}
\label{sec:top_wall_mid}

A subtler differentiator: fair-value estimation.  Most teams used
the raw mid of best bid/ask.  Frankfurt Hedgehogs used the
``Wall Mid'': the average of the bid-wall and ask-wall prices,
order-book levels with consistently deep market-maker-bot
liquidity.

The raw mid is distorted by aggressive takers overbidding or
undercutting.  The Wall Mid filters this by anchoring to the
deepest levels, corresponding to bots that ``appeared to know the
true price and simply quoted rounded versions of it ($\pm 2$
ticks).''

They validated empirically: buying or selling one lot on the
official site, PnL is computed against the true internal price.
Wall Mid tracked it to sub-tick accuracy; raw mid drifted by
several ticks in volatile periods.

This mattered most for Kelp (slow random walk).  Quoting around
the Wall Mid instead of raw mid reduced adverse-selection losses
on passive orders.  Frankfurt Hedgehogs estimate Kelp generated
$\sim$5{,}000 SeaShells/round, modest in absolute terms but
entirely from fair-value precision.

\subsection{Edge-case handling}
\label{sec:edge_cases}

Several technical details separated top implementations from
merely correct ones:

\begin{itemize}
  \item \textbf{Position limits.}  CMU Physics found fully hedging
        both Picnic Baskets simultaneously was impossible under
        constituent limits.  Solution: trade the \emph{difference}
        in basket premiums ($P_1 - P_2$), using 100\% of Basket~1's
        limit and 60\% of Basket~2's, then z-score the remainder
        and market-make with the last 8\%.

  \item \textbf{Integer rounding.}  Frankfurt Hedgehogs' Macaron
        sell orders priced at \texttt{int(externalBid + 0.5)}, the
        max price that still triggered fills from a hidden taker
        bot executing $\sim$60\% of eligible trades.  Wrong by
        1~SeaShell meant never filling (too high) or leaving edge
        (too low).

  \item \textbf{Memory and serialisation.}  CMU Physics' Round~3
        crash was Jasper's visualiser exceeding the 100\,MB Lambda
        limit.  Every byte of inter-timestep state matters.
        Frankfurt Hedgehogs stayed deliberately low-memory,
        storing only minimal rolling stats per product.

  \item \textbf{Conversion limits.}  Macarons had a 10-unit
        conversion limit per timestep.  Frankfurt Hedgehogs
        initially quoted 10 units/step, capturing $\sim$60\% of
        max profit.  Quoting 20--30 would have converted surplus
        inventory and pushed toward the optimal 130{,}000--160{,}000.
        CMU Physics independently found adjusting from 10 to 30
        Macarons/step nearly doubled volume (56k $\to$ 95k units).

  \item \textbf{Hedging cost analysis.}  CMU Physics tallied
        Volcanic Rock hedging trades and found 40{,}000+ SeaShells
        of spread cost.  Going unhedged (worst case $\sim$16{,}000)
        increased backtested PnL from 80k to 250k/day.

  \item \textbf{Storage fees and holding costs.}  Macarons had a
        0.1 SeaShell/step/unit storage fee.  Over 10{,}000 steps,
        holding the 75-unit limit would cost
        $75 \times 0.1 \times 10{,}000 = 75{,}000$ SeaShells, enough
        to wipe out most arbitrage profits.  CMU Physics sold
        locally and converted immediately, minimising holding.
        Teams accumulating inventory for better prices were eaten
        alive by storage.
\end{itemize}

\begin{keyfactbox}[Common patterns in top-10 edge-case handling]
The top teams share three engineering habits:
\begin{enumerate}
  \item \textbf{Compute all costs explicitly.}  Spreads, storage,
        conversion fees, tariffs are often larger than the gross
        edge.  Write them as formulas before coding.
  \item \textbf{Test against the memory limit.}  The $\sim$100\,MB
        Lambda limit is binding.  Verbose logging, large rolling
        windows, and visualiser libraries can silently push you
        over.
  \item \textbf{Use traderData serialisation defensively.}
        Cross-step state lives in \texttt{traderData}, but
        serialisation has overhead: test round-trip time and size
        before adding state.
\end{enumerate}
\end{keyfactbox}

\subsection{Manual-round preparation}
\label{sec:manual_prep}

Top teams approached manuals systematically, not intuitively.
Both Frankfurt Hedgehogs and CMU Physics modelled game-theoretic
manuals (Container, Suitcase) by computing Nash equilibria then
adjusting for behavioural priors.  CMU Physics built a probability
distribution over player types (Nash-followers, random, nice-number
biased) and Monte Carlo'd over it; Frankfurt Hedgehogs mapped P2
allocation patterns onto P3 containers.

For the Round~5 news-trading manual, Frankfurt Hedgehogs calibrated
expected price-move magnitudes from the identical P1/P2 manuals:
prior-year study again.

Key lesson: \emph{real players do not play Nash}.  Both teams
found their Nash calculations overestimated broad-population
rationality.  CMU Physics modelled five Container types:
Nash-followers (15\%), random (50\%), nice-number-biased (20\%),
naive-EV (10\%), flawed-MC (5\%).  Frankfurt Hedgehogs used P2$\to$P3
mapping for over/under-selection.  After Round~2 results, CMU
revised priors upward on Nash-following ($\sim$50--60\%) and
confirmed nice-number bias (37, 73).  This is Bayesian updating:
prior-year outcome as prior, current round as likelihood.

Further: Frankfurt Hedgehogs noted that opening a second Round~2
container (50{,}000 SeaShells) exceeded the Nash EV (expected
value, the probability-weighted average payoff) of any single
container, making the second pick negative-EV.  CMU Physics
reached the same conclusion by Monte Carlo.  Teams who opened two
anyway (``more picks $=$ more diversification'') lost money at
the margin.  Quantitative reasoning overrode intuition.


%======================================================================
\section{The hardcoding exploit}
\label{sec:hardcoding}
%======================================================================

In P3 Rounds~1--2, bot behaviour replicated earlier competitions
with 95\%+ accuracy.  True prices were regenerated (unlike P2,
which reused some price paths), but bot \emph{actions} (exact
timesteps, sizes, and directions of market orders) matched
historical data almost perfectly.

\subsection{What teams actually did}

Teams that noticed could ``hardcode'' algorithms to frontrun
predictable bot actions.  E.g.\ if a bot historically bought
Rainforest Resin at timestep~600, pre-buy all ask liquidity just
before and sell to the bot higher.  Some teams made millions of
SeaShells this way in Rounds~1--2.

Frankfurt Hedgehogs found this early and frontrun, citing the
unaddressed P2 precedent (Linear Utility's writeup).  They added
an automatic fallback: revert to the non-hardcoded strategy if
bot behaviour diverged from history.

CMU Physics learned the exploit the hard way.  After Round~1
they were 771st.  When IMC re-ran after ruling hardcoding as
cheating, they jumped to 9th, a 762-place swing.

\subsection{When hardcoding is rational and when it is a trap}

\begin{pitfallbox}[The hardcoding temptation]
Hardcoding tempts because it gives huge backtests with zero
parameter risk: the ``strategy'' is a lookup table.  It is a
trap for three reasons:
\begin{enumerate}
  \item \textbf{Fragility.}  If IMC changes bots (as after P3
        Round~2), hardcoded strategies produce zero PnL.
        All-in-on-hardcoding teams had no fallback.
  \item \textbf{Rule risk.}  IMC banned hardcoding and required
        re-submissions; caught teams dropped precipitously.
  \item \textbf{Opportunity cost.}  Reverse-engineering bot
        timesteps is time not spent on robust strategies that
        compound over five rounds.
\end{enumerate}
\end{pitfallbox}

The rational response (Frankfurt Hedgehogs) was to build
hardcoding \emph{with a fallback}: exploit while it works, revert
automatically when it breaks.  Same principle as production systems:
no single data source without a degraded-mode alternative.

After Round~2, Frankfurt Hedgehogs reported the exploit to IMC
even though fixing it cost their advantage.  IMC fixed bots for
Rounds~3--5 and allowed re-submissions.  Meta-lesson: in
competitions and industry, exploits based on implementation bugs
rather than genuine structure get patched.

\subsection{The P2 precedent and the ethics of exploits}

Hardcoding was not new to P3.  In P2, Linear Utility documented
that some price paths were directly reused, making \emph{prices}
(not just bot behaviour) predictable.  Their writeup became a
key P3 reference: Frankfurt Hedgehogs hardcoded in P3 precisely
because ``a similar situation had gone unaddressed in~P2.''

Frankfurt Hedgehogs' decision to report the exploit reflects a
principle: exploits riding bugs rather than genuine structure are
fragile, reputation-damaging, and value-destroying.  In a
competition, IMC patches and your edge vanishes; in industry, the
counterparty figures it out, the regulator investigates, or the
venue changes its matching rules.  Sustainable edge comes from
market structure, not implementation accidents.

Lesson for P4: assume IMC makes bot behaviour harder to predict
from prior-year data, but also assume \emph{some} predictable bot
behaviour persists: the competition needs bots to provide
liquidity and create opportunities, and fully random bots would kill
the alpha.  The skill is distinguishing structural patterns (which
persist) from implementation artefacts (which get patched).

Frankfurt Hedgehogs even tried reverse-engineering the random
seed, brute-forcing all 4 billion NumPy seeds on a Raspberry Pi
for 24 hours against the first $\sim$100 PnL returns.  No match,
but it shows the creative lengths top teams go to.

\begin{historybox}
Prosperity~1 launched in 2023 with Amethysts (fixed-value MM) and
Starfruit (random walk).  By P3 (2025), participation hit 12{,}000+
worldwide.  Format evolved: P2 added conversion (Orchids) and ETF
baskets; P3 added options (Volcanic Rock Vouchers), cross-island
arb (Macarons), and the controversial hardcoding exploit.  IMC
uses Prosperity as a recruiting pipeline: top finishers get
interview invitations.  P3 rules excluded prior top-10 teams from
prize money (Frankfurt Hedgehogs still competed for ranking).
\end{historybox}


%======================================================================
\section{What to watch for in P4}
\label{sec:p4_watch}
%======================================================================

Based on P1--P3 archetype evolution, some product-round patterns
are predictable, others deliberately novel.

\subsection{Expected archetype mapping}

The product--archetype mapping across all three completed
editions:

\begin{center}
\small
\begin{tabular}{@{}llll@{}}
\toprule
\textbf{Archetype} & \textbf{P1 (2023)} & \textbf{P2 (2024)}
  & \textbf{P3 (2025)} \\
\midrule
Fixed-value MM & Amethysts & Amethysts & Rainforest Resin \\
Drifting-value MM & Starfruit & Starfruit & Kelp \\
Volatile / chaotic & --- & --- & Squid Ink \\
Basket / ETF arb & --- & Gift Basket & Picnic Baskets 1\&2 \\
Options / derivatives & --- & Coconut options
  & Volcanic Rock Vouchers \\
Cross-exchange arb & --- & Orchids & Macarons \\
Informed bot & (unnamed) & (unnamed) & Olivia \\
\bottomrule
\end{tabular}
\end{center}

\noindent Each edition adds complexity to existing archetypes
(one basket $\to$ two; one option $\to$ five strikes) plus one
genuinely novel element (Squid Ink's volatility in P3, Orchids'
sunlight/humidity in P2).

Concrete P4 expectations:

\begin{itemize}
  \item \textbf{Round~1: flat-value market making.}  Every edition
        opened with a near-constant-value tutorial product
        (Amethysts, Rainforest Resin).

  \item \textbf{Rounds~2--3: basket stat-arb.}  Gift Basket (P2),
        two Picnic Baskets (P3).  P4 may add more baskets or
        time-varying compositions.

  \item \textbf{Rounds~3--4: derivatives.}  Coconut options (P2),
        Volcanic Rock Vouchers across five strikes (P3).  Puts or
        American-style plausibly next.

  \item \textbf{Round~4--5: cross-exchange / conversion arbitrage.}
        Orchids (P2), Macarons (P3).  Conversion mechanics grow
        more layered each year.

  \item \textbf{All rounds: an informed trader.}  Every edition
        has had at least one predictive bot trading daily extrema
        (Olivia in P3).  IMC may change the pattern; some form
        persists.
\end{itemize}

\begin{keyfactbox}[P4 preparation checklist]
\begin{enumerate}
  \item Build or fork a backtester and verify it works on P3 sample
        data before the competition starts.
  \item Build or install a visualisation dashboard with order-book
        display, PnL tracking, and price normalisation.
  \item Read all available P1--P3 writeups (Frankfurt Hedgehogs,
        CMU Physics, Linear Utility, Stanford Cardinal, jmerle,
        Alpha Animals at minimum).
  \item Prepare skeleton strategies for each archetype: fixed-value
        market making, drifting-value market making, basket
        stat-arb, option pricing / IV scalping, conversion
        arbitrage, and informed-bot detection.
  \item Join the Discord server on day one; moderator hints and
        community analysis are decisive information sources.
  \item Set up \texttt{traderData} serialisation and test it
        under the 100\,MB Lambda memory limit.
  \item Submit a working (even if basic) algorithm within hours
        of each round opening.
  \item After each round, re-optimise parameters on the latest
        data, choosing from plateau regions rather than peaks.
  \item Analyse all bot trades by quantity, timing, and direction
        before trader IDs are revealed.
  \item Designate team roles: one person owns manual rounds, one
        owns tooling, and at least one focuses on new-product
        analysis.
\end{enumerate}
\end{keyfactbox}

\subsection{What IMC might change}

P4 uses XIRECs and a deep-space theme, breaking from the
island/seashell motif.  Prior top-10 teams are now excluded from
prize rankings (still compete for ranking), suggesting IMC is
levelling for returning teams' structural advantage.  The
hardcoding exploit is patched, and bots will likely vary more
across rounds.  Novel mechanics (futures with margin, multi-leg
spreads, path-dependent payoffs) are plausible given participant
sophistication.

\subsection{Discord as an information source}

Frankfurt Hedgehogs call the Discord server ``extremely valuable''
for moderator tips and clarifications, but warn it is ``full of
noise'' from inexperienced participants plus occasional
psychological warfare: fake backtesting screenshots designed to
intimidate competitors into impulsive changes.

Recommended: monitor Discord for moderator hints and API/rule
clarifications; never share your strategy; never pivot on
unverified claims.  ``Stick to your own validation methods, and
avoid the trap of overfitting or changing your approach
impulsively just because of noise on Discord.''

\begin{pitfallbox}[Discord psychological warfare]
In P3, Discord posts showed extraordinary backtest numbers, some
from massively overfit strategies that would collapse live, some
outright fabrications.  If a competitor claims 500{,}000+ on one
product, do not panic; verify your own strategy against your own
data.  Teams that pivoted on Discord noise consistently
underperformed their own backtests.
\end{pitfallbox}

\subsection{The ``Wall Mid'' fair-value concept}

Frankfurt Hedgehogs' ``Wall Mid'' is the average of the bid-wall
and ask-wall prices, the levels with consistently deep designated-MM
liquidity, typically corresponding to the true internal price
$\pm 2$ ticks.  Averaging them gives a more stable, accurate fair
value than the raw mid (which is distorted by aggressive takers).

They verified by buying or selling one lot on the official site
and reading PnL vs.\ the true internal price.  This
ground-truth-checking methodology is worth replicating for any
Prosperity product with an official backtesting environment.

For P4, the specific Wall Mid form may change (different quoting
patterns, tick sizes), but the principle endures: true fair value
is rarely the simple best-bid/ask mid; it is usually better
estimated from deep-liquidity levels.

\subsection{Backtesting discipline}
\label{sec:backtest_discipline}

A final differentiator: top teams used backtesting as a
\emph{falsification} tool, not a confirmation tool.  Frankfurt
Hedgehogs warn: ``never optimize purely for website score.  Doing
so is extremely prone to overfitting on simulation-specific
randomness rather than building strategies that generalize.''
They used two complementary approaches:

\begin{itemize}
  \item \textbf{Jasper's backtester} for taking/simple-quoting
        strategies: fast, flexible, parameterised for systematic
        optimisation.
  \item \textbf{The official website} for bot-interaction
        strategies (Rainforest Resin MM, Macaron hidden-taker),
        since the local tool could not replicate bot fill
        probabilities.
\end{itemize}

Linear Utility's custom backtester output logs in the website's
exact format, enabling direct comparison.  Discrepancies were
treated as bugs, never noise.

CMU Physics learned this the hard way.  Their Squid Ink
spike-detector was backtested on limited data and happened to
work on the Round~1 re-run; they recognised luck, not edge.
They cut Squid Ink to 10\% of max position and focused on robust
signals, trading upside for stability.


%======================================================================
\section{What cannot be planned}
\label{sec:cannot_plan}
%======================================================================

Even the best-prepared teams encounter irreducible uncertainty.

\textbf{Novel products.}  Every edition has at least one product
that does not fit prior archetypes.  P3's Squid Ink, with extreme
volatility and unpredictable spikes, defeated standard
mean-reversion and market-making.  CMU Physics: ``Squid Ink was
pure chaos---with regular 100 SeaShell swings within a single step
and no obvious structure.''  When theory fails, reduce position
and gather live data rather than forcing a model.

\textbf{Regime changes mid-round.}  Frankfurt Hedgehogs' quadratic
IV fit for Volcanic Rock Vouchers worked in sample but ``stopped
being a good model on the submission day,'' producing losing
positions all day.  A rolling-window would have adapted; the
static fit could not.

\textbf{Infrastructure failures.}  CMU Physics' Round~3 crash
(visualiser memory) was unpredictable from local testing and
missed by the official backtester.  The only mitigation is
defensive programming: memory well below limits, worst-case-volume
testing, kill-switches for runaway state.

\textbf{Team coordination under pressure.}  Both teams describe
late-night sessions (``6~a.m.\ discussing Squid Ink dynamics'')
and fatigued judgment calls.  When CMU Physics dropped to 241st
after Round~3, they ``felt defeated and were honestly ready to
give up.''  They did not, and the next round was their best.
Composure after a setback is a character trait, not a strategy.

\textbf{Stochastic outcomes despite correct strategy.}  Frankfurt
Hedgehogs led every round except the last, where an unknown team
posted 850{,}000 SeaShells (vs.\ the typical 200--400k top range)
to take first.  They accepted it: ``in any probabilistic
environment \ldots\ some degree of randomness is inevitable.''
CMU Physics' Round~1 Squid Ink happened to spike \emph{for} them
on the re-run: pure luck in their 9th-place finish.

\textbf{The spread cost of hedging.}  CMU Physics found hedging
Volcanic Rock Vouchers cost 40{,}000+ SeaShells/day in spread: the
underlying had a 1-tick spread and frequent rebalancing piled up.
The ``correct'' textbook move (always hedge delta) was wrong in
this microstructure.  Frankfurt Hedgehogs concurred: ``explicit
delta hedging would have been prohibitively expensive given the
bid-ask spreads.''  Microstructure costs can invalidate
theoretically optimal strategies, and the true costs only appear
live.

\textbf{Teammate dynamics.}  CMU Physics had never spoken before
the competition; Frankfurt Hedgehogs had played P2 together.
Both configurations worked, but both teams describe tension:
disagreements, sleep deprivation, the weight of rank drops.
No formula for team chemistry, and no backtester can simulate
whether a teammate will push through a bad round.

\begin{intuitionbox}
The unplannable share a structure: they are \emph{out-of-sample}
(OOS, outside the historical data used to build and test the
strategy) risks no in-sample analysis can eliminate.  Novel products are
OOS by definition; regime changes make your calibration
non-representative; infrastructure failures sit outside the tested
distribution; team dynamics have no historical data.  Response:
maintain margin.  Keep positions small enough that one surprise
doesn't eliminate you, code modular enough to swap strategies
fast, team cohesive enough that a bad round is data, not crisis.
\end{intuitionbox}


%======================================================================

\begin{prosperitybox}
\textbf{Meta-lesson.}  The best predictor of Prosperity success
is not mathematical sophistication but \emph{operational
readiness}.  Teams with tooling, skeletons, and prior-year
knowledge on day~1 beat teams with deeper theory who spent two
days building infrastructure.  Frankfurt Hedgehogs' ``weeks of
careful preparation'' and Linear Utility's ``in-house tools''
let both teams focus on strategy from timestep~1, while most
competitors were still parsing API docs.
\end{prosperitybox}

\begin{gsbox}
The same lesson applies on a Goldman Sachs desk at higher stakes.
The best strat is not the one with the most elegant mathematics
but the one in production, monitored, and iterable.  A model
taking three months to deploy is dominated by a simpler one that
ships in a week and refines over the next eleven.  Desk traders
value reliability and responsiveness over theoretical optimality.
An off-cycle engineer who ships fast, iterates on live feedback,
and handles edge cases earns trust, the same trust Prosperity
teams earn by submitting early rather than optimising in
isolation.
\end{gsbox}


%======================================================================
\section{Key facts to memorise}
\label{sec:top_team_key_facts}
%======================================================================

\begin{enumerate}
  \item \textbf{Tooling before strategy.}  Build backtester,
        visualiser, and analysis notebooks before the competition
        starts.  Use prior-year data to validate the pipeline.

  \item \textbf{Study all available writeups.}  Product archetypes
        repeat across Prosperity editions.  Detection of the
        Olivia bot in P3 was directly enabled by reading P1 and
        P2 writeups.

  \item \textbf{Simple beats complex.}  The winning formula is:
        fair-value estimate $+$ fixed threshold $+$ inventory
        management.  Every top team says this; none deployed
        complex ML in their final submission.

  \item \textbf{Submit early, iterate live.}  Live-data feedback
        is worth more than another day of backtesting.  CMU
        Physics' comeback from 241st to 8th in one round
        demonstrates the power of rapid iteration.

  \item \textbf{Detect bot patterns.}  Filter trades by quantity,
        timing, and direction.  The informed-bot signal has been
        the single highest-value alpha source in every Prosperity
        edition.

  \item \textbf{Reuse signals across products.}  Olivia's
        Croissants position biased basket spread thresholds.
        Cross-product signal integration compounded edges for
        both Frankfurt Hedgehogs and CMU Physics.

  \item \textbf{Choose parameter plateaus, not peaks.}  Select
        parameters from flat regions of the PnL landscape.
        Sensitivity to small parameter changes is the primary
        warning sign of overfitting.

  \item \textbf{Quantify worst-case losses before deploying.}
        CMU Physics' unhedged Volcanic Rock position and Frankfurt
        Hedgehogs' VaR-bounded mean-reversion exposure both
        succeeded because the teams computed realistic downsides
        first.

  \item \textbf{Handle edge cases explicitly.}  Position limits,
        integer rounding, memory limits, conversion caps, and
        serialisation overhead have all cost top teams tens of
        thousands of SeaShells when handled carelessly.

  \item \textbf{Operational readiness beats mathematical
        sophistication.}  The team that ships a working model on
        day one and iterates will beat the team that ships a
        perfect model on day three.
\end{enumerate}
